\section{Einleitung} \label{sec:einleitung}

\subsection{Problemstellung} \label{sec:problemstellung}

Die radiologische Befundung gehört zu den ressourcenintensivsten Aufgaben im klinischen
Alltag. Röntgenaufnahmen, CT- und MRT-Untersuchungen müssen von hochspezialisierten
Fachärztinnen und Fachärzten analysiert und befundet werden -- ein Prozess, der weder
einfach skalierbar noch uneingeschränkt fehlertolerant ist. Das Untersuchungsvolumen
ist in den vergangenen zehn Jahren um drei bis fünf Prozent jährlich gewachsen
\cite{aiSecondReader2024}, während die Fachkräftebasis nicht proportional mitgewachsen
ist: In Deutschland sind von 13.791 approbierten Fachärztinnen und Fachärzten für
Radiologie lediglich 10.139 berufstätig \cite{statistaRadiologen2024}.

Die Folgen dieser Überlastung sind messbar: Die normierte Arbeitslast zum Zeitpunkt
diagnostischer Fehler im CT-Befund betrug im Schnitt 121~\% des Normalwerts
\cite{workloadErrors2023}; 44~\% der männlichen und 65~\% der weiblichen Radiologen
berichten Burnout-Symptome \cite{burnoutACR2024}. Die Diagnoseleistung ist damit
strukturell von Faktoren wie Arbeitsbelastung und Ermüdung abhängig.

\subsection{Motivation: Methodischer Transfer aus der industriellen Bildanalyse}
\label{sec:motivation}

In der industriellen Qualitätskontrolle sind KI-gestützte Bildanalysesysteme seit über
einem Jahrzehnt produktiv im Einsatz. Die methodische Grundlage dieser Systeme --
Convolutional Neural Networks, Transfer Learning, Bildsegmentierung und Klassifikation --
beschreibt allgemeine Verfahren zur Anomalieerkennung in Bilddaten, die auf die
medizinische Bildgebung übertragbar sind. Empirische Studien belegen, dass der Transfer
in die medizinische Diagnostik zu Ergebnissen auf radiologischem Niveau führt
\cite{chexnet2017}. Dieser Artikel untersucht die methodischen Grundlagen, Leistungsmetriken
und Systemarchitekturen, die diesen Transfer ermöglichen.

\subsection{Forschungsfragen} \label{sec:zielsetzung}

Der Artikel richtet sich an Studierende der Gesundheitstechnik mit technischem
Grundverständnis und untersucht drei Forschungsfragen:

\begin{enumerate}
    \item Welche technischen Verfahren der industriellen Bildanalyse sind auf die
          medizinische Diagnostik übertragbar?
    \item Wie lässt sich der Nutzen KI-gestützter Systeme quantitativ messen und belegen?
    \item Welche Systemarchitektur ermöglicht eine DSGVO-konforme Implementierung
          im Krankenhaus?
\end{enumerate}
